\section{Espaços topologicos}

\begin{definition}
    Chamamos de \emph{topologia} em \(X\) uma familia \(\tau \subset \wp(X) :\) (T1) \(\varnothing, X \in \tau\); (T2) \(\mathcal{C} \subset \tau \Rightarrow  \bigcup \{U : U \in \mathcal{C}\} \in \tau\) e (T3) \ \( U,  V \in \tau\Rightarrow  U\cap V \in \tau \). 
\end{definition}

\begin{note}
    Um subconjunto \(F\ce\subset X \ \leftrightharpoons \ (X \setminus F)\ab \subset X\), como no {\scriptsize \(\boxtimes\)} \hyperlink{coro27}{2.7.}, existe uma versão equivalente da definição topologia em termos de fechados.
\end{note}

\begin{definition}
    Dadas duas topologias \(\tau_1, \ \tau_2 \) em \(X\) tais que \(\tau_1\subset \tau_2\). Então, dizemos que \(\tau_1\) é mais \emph{fraca} ou \(\tau_2\) mais \emph{forte}  que a outra.  
\end{definition}

\ejemplo{
    \begin{itemize}[left =0cm, label = {\tiny\(\bullet \)}]
    \item[] \- \hspace{-0.6cm}\textcolor{verdeoscuro}{\underline{\textbf{Exemplos}}} \hspace{1cm}{\tiny\(\bullet \)} \ \((X, d)\) com \(\tau = \tau_d \) (esp. métricos) \hfill \-  
    \item \(\tau = \{\varnothing, X\}\) (top. trivial) \hfill \(\mid\) \hfill {\tiny\(\bullet \)} \ \(\tau = \{\varnothing, \{a\}, \{a,b\}\}\) (top. de Sierpinski)
    \item \(\overline{\tau} = \{F \subset X : F \text{ finito}\}\) (top. cofinita)\hfill \(\mid\) \hfill {\tiny\(\bullet \)} \ \(\tau = \wp(X)\) (top. discreta)  
\end{itemize}
}

\subsection*{Aderência e interior}

\begin{definition}
    Dados \((X,\tau)\) e \(A\subset X\), definimos \emph{aderência} e \emph{interior} de \(A\), respectivamente, 
    \[\overline{A}:= \bigcap \{F\ce \subset X : A \subset F \} \text{ \ \ e \ \ } A^\circ := \bigcup \{U\ab \subset X : U \subset A\}.\]  
\end{definition}
\vspace{-0.2cm}
\hypertarget{prop4x}{}
\begin{proposition}
    \textbf{4.*.} As aplicações \(\operatorname{f}:  A \mapsto \overline{A}   \) e \(\operatorname{i} : A \mapsto A^\circ\) verificam: 
    \vspace{-0.2cm}
    \begin{enumerate}[label = (\(\operatorname{f}\)\arabic*), left = 0.75cm]
        \item \(A\subset \overline{A}\). \hfill (\(\operatorname{f}\)2) \ \(\overline{A\cup B} = \overline{A} \cup \overline{B}\).\hfill (\(\operatorname{f}\)3) \ \(A\ce\subset X \ \leftrightharpoons \ A = \overline{A}\)\hspace{0.1cm}.
        \item[(\(\operatorname{i}\)1)] \(A^\circ \subset A\). \hspace{1.65cm} (\(\operatorname{i}\)2) \ \((A \cap B)^\circ = A^\circ \cap B^\circ\).\hspace{1.1cm} (\(\operatorname{i}\)3) \ \(A\ab \subset X \ \leftrightharpoons \ A = A^\circ\).
    \end{enumerate}
\end{proposition}
\begin{note}
    Tem versões mais completas da \(\square\) \hyperlink{prop4x}{4.*.}, até da pra definir topologia a partir de aplicações verificando aquelas propriedades, olha se queser \cite[Pág. 9]{mujica2005}.  
\end{note}
\begin{proposition}
    \textbf{4.5.} \(X\setminus \overline{A} = (X \setminus A)^\circ\) e \( X\setminus A^\circ = \overline{X\setminus A}\). \comentario{Mais uma vez Ex. \hyperlink{demorgan}{1.A.}}
\end{proposition}

\subsection*{Sistemas de vizinhanças}

\begin{definition}
    \(U\subset X\) é \emph{vizinhança} de \(x \in X\) se \(x \in U^\circ\), cujo casso \( U \in \U_x\). 
\end{definition}

\hypertarget{prop52}{}
\begin{proposition}
    \textbf{5.2.} As \emph{vizinhanças} \(U \in \U_x\) verificam as seguintes propriedades: 
\vspace{-0.2cm}
    \begin{enumerate}[label = (\alph*), left = 0.6cm]
        \item \(x \in U\). \hfill  (b) \ \( V \in \U_x \Rightarrow U\cap V \in \U_x\). \hfill (d) \ \(U \subset V\subset X \Rightarrow V \in \U_x\).
        \item[(c)] \(\exists V \in \U_x : V\subset U\) e \(\forall y \in V, \ U \in \U_y\). \hfill (e) \  \(U\ab\subset X \ \leftrightharpoons \ \forall x \in U, \  U \in \U_x\). 
    \end{enumerate}
\end{proposition}
\begin{note}
    Mais uma vez, é possível definir a topologia \(\tau = \{U \subset X : U \in \U_x, \ \forall x \in U\}\) a partir de uma familia verificando tais propriedades, \cite[Pág. 12]{mujica2005}.
\end{note}

\begin{definition}
    Uma familia \(\mathcal{B}_x\subset \U_x\) é \emph{base de vizinhanças} se \(\forall U \in \U_x, \ \exists V \in \mathcal{B}_x : V \subset U\). 
\end{definition}

\begin{proposition}
    \textbf{5.7.} Pra cada \(x \in X \) seja \(\mathcal{B}_x\) familia verificando: 
    \vspace{-0.2cm}
    \begin{enumerate}[label = (\alph*), left = 0.6cm ]
        \item \(\forall U \in \B_x, \ x \in U\). \hfill (b) Dados \(U,V \in \B_x, \ \exists W \in \B_x : W \subset U \cap V\). 
        \item[(c)] Dada \(U \in \B_x, \ \exists V \in \B_x, \ V\subset U : \forall y \in V, \ \exists W \in \B_y \) tal que \( W \subset U\). 
    \end{enumerate}
    Então, \(\tau := \{U \subset X : \forall x \in U, \ \exists V \in \B_x \mid V \subset U\}\) é uma topologia em \(X\) e \(\B_x\) uma base de vizinhanças pra \(\tau\). \comentario{Considere \(\U_x := \{U\subset X : \exists V \in \B_x \mid V \subset U\}\), a afirmação segue da recíproca da \(\square\) \hyperlink{prop52}{5.2.} comentada previamente}
\end{proposition}

\begin{proposition}
    \textbf{5.8.} Sejam \(A \subset X\) e \(\B_x\) base de vizinhanças. Então, \(\overline{A} = \{x \in X : \forall V \in \B_x, \ V\cap A \neq \varnothing\}\) \ e \ \(A^\circ = \{x \in X : \exists  V \in \B_x : V\subset A\}\).
\end{proposition}

\begin{example}
    Em \((X, d)\), as bolas (abertas ou fechadas) são bases de vizinhanças. 
\end{example}
\begin{definition}
    \((X, \tau)\) satisfaz o \emph{primeiro axioma de enumerabilidade} se cada \( x \in X\) admite base de vizinhaças \(\B_x\) enumerável.  
\end{definition}

\subsection*{Bases para os abertos}

\begin{definition}
    Diz-se que \(\B \subset \tau\) é \emph{base} de \(\tau \) se \(\forall U \ab \subset X, \ \exists \mathcal{C} \subset \B : U = \bigcup \{V : V \in \mathcal{C}\}\). 
\end{definition}

\begin{proposition}
    \textbf{6.3.} \(\B \subset \tau\) é base \(\leftrightharpoons \ \forall x \in U\ab\subset X, \  \exists V \in \B : x \in V \subset U\), pra cada \( U \in \tau \). 
\end{proposition}

\hypertarget{prop66}{}
\begin{proposition}
    \textbf{6.6.} Seja \(\B \subset \wp(X)\) verificando: (a) \(X = \bigcup\{V : V \in \B\}\) e (b) Dados \(U, V \in \B \text{ e } x \in U \cap V, \ \exists W \in \B : x \in W \subset U \cap V \). Então, os conjuntos da forma:  
    \[U = \bigcup \{V : V \in \mathcal{C}\} \text{ \ \ com \ \ } \mathcal{C} \subset \B,\]
    definem uma topologia \(\tau\) em \( X\), pra a qual \(\B\) é base. 
\end{proposition}

\ejemplo{%\centering \(+\) Exemplos importantes (em \(X = \R\))
\begin{example}
    {\scriptsize \(\bullet\)} \ \(\B = \{(a, \infty) : a \in \R\}\) \hfill \(\mid \) \hfill{\scriptsize \(\bullet\)} \ \(\B :=\{[a,b): a <b\}\) (top. Sorgenfrey) 
\end{example}
}

\begin{definition}
    \((X, \tau)\) satisfaz o \emph{segundo axioma de enumerabilidade} se \(\exists \B\subset \tau\) base enumerável.  
\end{definition}

\begin{example}
    Em \((\R, \tau_2)\), a familia \(\B := \{(a,b) : a < b\}\) é base de \(\tau_2\), enumerável se tomarmos só \(a,b \in \mathbb{Q}\).  
\end{example}



\subsection*{Subespaços}

\begin{definition}
    Dados \((X,\tau )\) e \(S\subset X\), é claro que \(\tau_S = \{U \cap S : U\ab \subset X\}\) é uma topologia em \(S\) que chamamos de \emph{topologia induzida}, denotamos \(S\leq X\). 
\end{definition}

\begin{proposition}
    \textbf{7.3.} (a) \ \(U\ab\subset X \ \leftrightharpoons \ \exists U_1\ab \subset X : U = U_1 \cap S\). 
    \vspace{-0.25cm}
    \begin{enumerate}[label = (\alph*), left = 0.6cm]
        \item[(b)] \(F\ce \subset S \ \leftrightharpoons \ \exists F_1\ce \subset X : F = F_1 \cap S\). 
        \item[(c)] Se \(A \subset S\), então \(\overline{A}^S = S \cap \overline{A}^X\). 
        \item[(d)] \(U \in \U_x\) (em \(S\)) \(\leftrightharpoons\) \(\exists U_1 \in \U_x\) (em \(X\)) \( : U = U_1 \cap S\). 
    \end{enumerate}
\end{proposition}
